\Article{1}
At any place on Earth, the force that gives a definite direction to a magnetic needle suspended at its center of gravity is called the Earth's magnetic force, insofar as the seat of its cause can be sought only in the body of the Earth itself. This assumes that every extraneous cause capable of acting on the needle has been removed, such as a nearby artificial magnet or a nearby conductor carrying a galvanic current. It is indeed an open question whether the immediate causes of the regular and irregular hourly changes in this force might lie outside the body of the Earth. We may hope that the general attention natural scientists are now giving to these phenomena will yield significant insight in the future. But we must remember that these changes are comparatively very small. There must therefore be a much stronger, persistently acting principal force, which we take to have its seat in the Earth itself. It follows at once that the observational basis for investigating this principal force ought properly to be freed from those anomalous changes. This can only be done by taking \SourceLeaf{130}averages from numerous sustained observations. As long as such results, free of these changes, are unavailable for a large number of points across the entire surface of the Earth, the best that can be achieved is an approximation in which differences of the order of these anomalies may remain.

\Article{2}
The premise of our investigation is that the Earth's magnetic force is the combined action of the magnetized parts of the Earth's body. We picture magnetization as a separation of the magnetic fluid. Once we adopt this picture, the way these fluids act belongs among established physical facts. Like kinds of fluid repel and unlike kinds attract, inversely as the square of the distance. Replacing this picture with \textsc{Ampère}'s, which dispenses with magnetic fluids and holds that magnetism consists solely in persistent galvanic currents in the smallest parts of bodies, would change nothing in the results. The same would hold if we attributed terrestrial magnetism to a mixed cause, arising partly from separation of the magnetic fluids within the Earth and partly from galvanic currents within it. As is well known, for any galvanic current we can substitute a particular distribution of magnetic fluids on a surface bounded by the line of the current, such that it exerts exactly the same magnetic action as the galvanic current at every point of the space outside.

\Article{3}
To measure the magnetic fluids, we take as the positive unit, as in \textit{Intensitas vis magneticae} etc., that quantity of northern fluid which, at the distance chosen as the unit, exerts a motive force equal to the chosen unit of force on an equal quantity of the same fluid. When we speak simply of the magnetic force at any point in space due to magnetic fluid elsewhere, we always mean \SourceLeaf{131}the motive force exerted there on a unit of positive magnetic fluid. In this sense, therefore, the magnetic fluid $\mu$, thought of as concentrated at a single point, exerts at the distance $\rho$ the magnetic force $\frac{\mu}{\rho\rho}$, repelling or attracting along the straight line $\rho$ according as $\mu$ is positive or negative. Let $a,b,c$ denote the coordinates of $\mu$ with respect to three axes intersecting one another at right angles, and $x,y,z$ the coordinates of the point at which the force is exerted, so that $\rho=\BareRoot{((x-a)^2+(y-b)^2+(z-c)^2)}$. If we resolve the force parallel to the coordinate axes, its components are
\[
\frac{\mu(x-a)}{\rho^3},\qquad \frac{\mu(y-b)}{\rho^3},\qquad \frac{\mu(z-c)}{\rho^3}
\]
which, as is easily seen, equal the partial derivatives of $-\frac{\mu}{\rho}$ with respect to $x,y$ and $z$.

If, in addition to $\mu$, other portions of magnetic fluid, $\mu',\mu'',\mu'''$ and so on, act while concentrated at points whose distances from the point of action are respectively $\rho',\rho'',\rho'''$ and so on, then the components of the total magnetic force that results, parallel to the coordinate axes, equal the partial derivatives of
\[
-\left(\frac{\mu}{\rho}+\frac{\mu'}{\rho'}+\frac{\mu''}{\rho''}+\frac{\mu'''}{\rho'''}+\text{etc.}\right)
\]
with respect to $x,y$ and $z$.

